\subsection{交换群中的可和族}

在本节中我们只讨论其合成法则写作加法的豪斯多夫交换拓扑群。只有最重要的结果才改写成乘法记号。

设 \(\mathbf{G}\) 是豪斯多夫交换群，\(\mathbf{I}\) 是任意指标集，\((x_{i})_{i\in \mathbf{I}}\) 是 \(\mathbf{G}\) 的点族，以 \(\mathbf{I}\) 为指标。对每个有限子集 \(\mathbf{J}\)（\(\mathbf{I}\) 的），我们让元素 \(s_{\mathbf{J}} = \sum_{i\in \mathbf{J}}x_i\)（\(\mathbf{G}\) 中的元素）与之对应，

并称它为族 \((x_{i})_{i\in I}\) 的对应于集合 \(J\) 的有限部分和。若用 \(\mathfrak{F}(I)\) 表示 \(I\) 的有限子集全体，我们就得到映射 \(J\to s_J\)（\(\mathfrak{F}(I)\) 到 \(G\) 中的映射）。现在，\(\mathfrak{F}(I)\) 关于关系 \(\subset\) 是一个有向集；因为若 \(J\) 与 \(J'\) 是 \(\mathfrak{F}(I)\) 的两个元素，则 \(J\subset J\cup J'\) 且 \(J'\subset J\cup J'\)，而 \(J\cup J'\) 是 \(I\) 的有限子集。设 \(\Phi\) 是有向集 \(\mathfrak{F}(I)\) 的截面滤子（第一章，§ 6，第 3 节）。

定义 1. 设 \((x_{i})_{i\in I}\) 是豪斯多夫交换群 G 的点族；设 \(\mathfrak{F}(I)\) 是指标集 \(I\) 的有限子集全体，并且对每个有限子集 \(J\)（\(I\) 的子集），设 \(s_J\) 是那些 \(x_{i}\)（满足 \(i\in J\) 者）的和。族 \((x_{i})_{i\in I}\) 称为可和的，如果映射 \(J\to s_J\) 关于截面滤子 \(\Phi\)（集合 \(\mathfrak{F}(I)\) 关于关系 \(\subset\) 定向后的截面滤子）有极限；这个极限称为族 \((x_{i})_{i\in I}\) 的和，记作 \(\sum_{i\in I}x_{i}\)（当不致引起歧义时，简记为 \(\sum x_{i}\)，甚至 \(\sum x_{i}\)）。

定义 1 等价于下述说法：族 \((x_{i})\) 可和且其和为 \(s\)，是指对每个邻域 \(\mathbf{V}\)（\(\mathbf{G}\) 中原点的邻域），存在有限子集 \(\mathbf{J}_0\)（\(\mathbf{I}\) 的），使得对每个有限子集 \(\mathbf{J} \supset \mathbf{J}_0\)（\(\mathbf{I}\) 中的），都有 \(s_{\mathbf{J}} \in s + \mathbf{V}\)。

若 G 写作乘法，并且对 I 的每个有限子集 J 令 \(p_{J} = \prod_{i \in J} x_i\)，则称族 \((x_i)\) 是可乘的，当映射 \(J \to p_J\) 关于滤子 \(\Phi\) 有极限时；这个极限称为族 \((x_i)\) 的积，记作 \(\prod_{i \in I} x_i\)。

注. 1) 当 I 是有限集时，定义 1 就化作有限族的和的通常定义。更一般地，若 I 是任意的且都有 \(x_{i} = 0\)，但指标 \(i\) 属于 I 的某个有限子集 J 者除外，则和 \(\sum_{i \in I} x_{i}\) 等于 \(\sum_{i \in J} x_{i}\)。

\begin{enumerate}
\def\labelenumi{\arabic{enumi})}
\setcounter{enumi}{1}
\item
  可和族的定义不涉及指标集 I 的任何次序，因此我们可以说这样定义的和的概念是交换的。更精确地说，我们有如下性质：设 \((x_{t})_{t\in \mathbf{I}}\) 是可和族，\(\varphi\) 是某指标集 K 到集合 I 上的双射；那么若令 \(y_{x} = x_{\varphi (x)}\)，则族 \((y_{x})_{x\in \mathbb{K}}\) 是可和的且与 \((x_{t})\) 有相同的和。因为若 \(s = \sum_{i\in I}x_i\)，且 \(\sum_{i\in J}x_i\in s + V\) 对包含 \(J_0\) 的每个有限子集 J 成立，那么 \(\sum_{x\in L}y_x\in s + V\) 对 K 的包含 \(\varphi^{-1}(J_0)\) 的每个有限子集 L 成立。
\item
  定义 1 还更一般地适用于豪斯多夫拓扑空间 X 中的点族，只要在 X 上已定义了一个写作加法的结合且交换的合成法则；因为定义 1 并不使用拓扑群的公理。
\end{enumerate}

